% =====================================================
% part7_svd.tex ― 第7部:行列分解と応用解析
% =====================================================
\part{行列分解と応用解析}

% =====================================================
\chapter{最小二乗法 ― 解けない方程式に最善の答えを}\label{chap:lsq}
% =====================================================

\begin{goal}
\begin{itemize}
\item 解のない $A\vect{x} = \vect{b}$ に対し、残差のノルム
$\|\vect{b} - A\vect{x}\|$ を最小にする\textbf{最小二乗解}を、
\textbf{正規方程式} ${}^t\!AA\vect{x} = {}^t\!A\vect{b}$ で求められるようになる。
\item 正規方程式が「$\vect{b}$ を $\mathrm{Im}\,A$ へ正射影せよ」という
第V部の幾何から出てくることを、証明つきで説明できるようになる。
\item 回帰直線や多項式の当てはめを、最小二乗問題として定式化して
解けるようになる。
\item 「予測値を作る」操作そのものを一つの行列(正射影行列
$P = A({}^t\!AA)^{-1}{}^t\!A$)にまとめ、その性質
($P^2 = P$, ${}^t\!P = P$)を使えるようになる。
\end{itemize}
\end{goal}

\section{動機 ― 解のない方程式と向き合う}

第II部で「解の存在しない連立方程式」を学んだとき、
われわれはそこで引き返した。しかし現実のデータ解析では、
\textbf{解がないのが普通}である。たとえば $100$ 個の観測データに
$2$ 個のパラメータの直線を当てはめれば、式 $100$ 本・未知数 $2$ 個の
方程式 $A\vect{x} = \vect{b}$ になり、観測誤差のせいで厳密な解はまず存在しない。
引き返す代わりに、問いを変える:

\begin{itembox}[l]{最小二乗問題}
$A\vect{x} = \vect{b}$ が解をもたないとき、誤差(残差)のノルム
\[
\| \vect{b} - A\vect{x} \|
\]
を\textbf{最小にする} $\vect{x}$ を求めよ。
\end{itembox}

「厳密に成り立たせる」ことをあきらめ、「破れ(残差)をできるだけ
小さくする」ことに目標を変えるのである。残差の\textbf{ノルム}を
最小化するとは、残差の各成分の\textbf{二乗和}を最小化することにほかならない
($\|\vect{r}\|^2 = r_1^2 + \cdots + r_m^2$)。「最小二乗法」の
名の由来である。この最小化を、微分に頼らず\textbf{幾何だけで}
解いてしまえるのが、第V部で積み上げた内積の理論の力である。

\section{正規方程式 ― 正射影の出番}

$A\vect{x}$ の動ける範囲は、$A$ の列ベクトルが張る部分空間
$W = \mathrm{Im}\,A$ である。よって問題は
「$W$ の中で $\vect{b}$ に最も近い点を探せ」となり、
第V部の深掘りで証明済みの事実――\textbf{最も近い点は正射影である}――が
そのまま答えを与える。最良の $A\vect{x}$ は $\vect{b}$ の $W$ への正射影であり、
そのとき残差 $\vect{b} - A\vect{x}$ は $W$ と直交する。
直交条件を $A$ の各列との内積 $= 0$ として書くと
${}^t\!A(\vect{b} - A\vect{x}) = \vect{0}$、すなわち:

\begin{teiri}[正規方程式]\label{thm:lsq-normal}
$m \times n$ 行列 $A$ と $\vect{b} \in \R^m$ に対し、
$\vect{x} \in \R^n$ について次は同値である。
\begin{enumerate}
\item $\vect{x}$ は $\|\vect{b} - A\vect{x}\|$ を最小にする(\textbf{最小二乗解})。
\item ${}^t\!A A\,\vect{x} = {}^t\!A\,\vect{b}$(\textbf{正規方程式})。
\end{enumerate}
正規方程式は($A\vect{x} = \vect{b}$ が解けなくても)\textbf{必ず解をもつ}。
さらに $A$ の列が線形独立なら ${}^t\!AA$ は正則で、最小二乗解はただ一つ
$\vect{x} = ({}^t\!AA)^{-1}\,{}^t\!A\,\vect{b}$ に定まる。
\end{teiri}

\noindent\textbf{(証明)}\quad
$W = \mathrm{Im}\,A$ とおき、$\vect{b}$ の直交分解
(定理 \ref{thm:gs-decomp})を $\vect{b} = \vect{b}_W + \vect{b}_{\perp}$
($\vect{b}_W \in W$, $\vect{b}_{\perp} \in W^{\perp}$)とする。

\textbf{(1 $\iff$ $A\vect{x} = \vect{b}_W$)}\
任意の $\vect{x}$ に対し $A\vect{x} \in W$ で、
$\vect{b} - A\vect{x} = \vect{b}_{\perp} + (\vect{b}_W - A\vect{x})$ は
「$W^{\perp}$ の元 $+$ $W$ の元」の和だから、ピタゴラスの定理
(定理 \ref{thm:inner-pyth})より
\[
\|\vect{b} - A\vect{x}\|^2
= \|\vect{b}_{\perp}\|^2 + \|\vect{b}_W - A\vect{x}\|^2
\ \ge\ \|\vect{b}_{\perp}\|^2.
\]
等号は $A\vect{x} = \vect{b}_W$ のとき、かつそのときに限る。
$\vect{b}_W \in W = \mathrm{Im}\,A$ だから、これを満たす $\vect{x}$ は
\textbf{必ず存在}する。

\textbf{($A\vect{x} = \vect{b}_W$ $\iff$ 2)}\
$A\vect{x} = \vect{b}_W$ なら $\vect{b} - A\vect{x} = \vect{b}_{\perp} \in W^{\perp}$。
逆に $\vect{b} - A\vect{x} \in W^{\perp}$ なら、
$\vect{b} = A\vect{x} + (\vect{b} - A\vect{x})$ が
「$W$ の元 $+$ $W^{\perp}$ の元」の分解になるので、分解の\textbf{一意性}
(定理 \ref{thm:gs-decomp})より $A\vect{x} = \vect{b}_W$。
そして「$\vect{b} - A\vect{x} \in W^{\perp}$」は、$W$ の生成系である
$A$ の各列 $\vect{a}_1, \dots, \vect{a}_n$ と直交すること、すなわち
${}^t\vect{a}_i(\vect{b} - A\vect{x}) = 0$($i = 1, \dots, n$)と同値であり、
$n$ 本まとめて書けば ${}^t\!A(\vect{b} - A\vect{x}) = \vect{0}$、
移項して ${}^t\!AA\vect{x} = {}^t\!A\vect{b}$ である。

\textbf{(列が線形独立なら ${}^t\!AA$ は正則)}\
${}^t\!AA\vect{x} = \vect{0}$ とすると、左から ${}^t\vect{x}$ を掛けて
\[
0 = {}^t\vect{x}\,{}^t\!AA\vect{x} = {}^t(A\vect{x})(A\vect{x}) = \|A\vect{x}\|^2
\]
より $A\vect{x} = \vect{0}$。列が線形独立だから $\vect{x} = \vect{0}$。
すなわち $n$ 次正方行列 ${}^t\!AA$ は $\mathrm{Ker} = \{\vect{0}\}$ をもち、
正則である(第III部)。 $\blacksquare$

解けない方程式 $A\vect{x} = \vect{b}$ の両辺に ${}^t\!A$ を掛けるだけで、
\textbf{必ず解ける}方程式に化ける――驚くほど実用的な定理である。

\begin{chui}
${}^t\!A$ を掛けるという操作は同値変形ではない(第I部で警告した
「両辺に行列を掛けると解が増えうる」の典型例である)。増えた解は
「元の方程式の解」ではなく「残差最小の解」だ、というのが
定理 \ref{thm:lsq-normal} の主張である。なお、$A\vect{x} = \vect{b}$ が
運よく解をもつ場合は残差 $\vect{0}$ が最小だから、
\textbf{最小二乗解 $=$ 通常の解}となり、矛盾は起きない(問 \ref{q:lsq-ex4})。
また、列が線形従属な場合は最小二乗解が無数に現れる
(反例 \ref{hanrei:lsq-dependent})。その場合にどれを選ぶかは、
第 \ref{chap:pinv} 章(擬似逆行列)で決着する。
\end{chui}

\begin{rei}[回帰直線]\label{rei:lsq-regression}
データ $(0, 1), (1, 2), (2, 2), (3, 3)$ に直線 $y = a + bx$ を当てはめる。
各点を代入した $4$ 本の式を行列で書くと
\[
A = \mat{1 & 0 \\ 1 & 1 \\ 1 & 2 \\ 1 & 3}, \qquad
\vect{b} = \mat{1 \\ 2 \\ 2 \\ 3}, \qquad
A\mat{a \\ b} = \vect{b} \ (\text{解なし}).
\]
正規方程式は
\[
{}^t\!AA = \mat{4 & 6 \\ 6 & 14}, \quad
{}^t\!A\vect{b} = \mat{8 \\ 15}, \quad
\mat{4 & 6 \\ 6 & 14}\mat{a \\ b} = \mat{8 \\ 15}
\]
となり、解いて $a = 1.1$, $b = 0.6$。回帰直線は $y = 1.1 + 0.6x$ である。
統計学の教科書に載っている回帰係数の公式は、
この正規方程式を解いた結果にほかならない。
\end{rei}

\begin{itembox}[l]{深掘り:回帰分析は幾何学である}
最小二乗法による回帰(OLS)は計量経済学・統計学の中心的手法だが、
その本性は徹頭徹尾、内積空間の幾何である。
残差ベクトルは説明変数の張る空間と\textbf{直交}する
(だから残差と説明変数の標本相関は厳密に $0$ になる)。
当てはまりの良さを測る決定係数 $R^2$ は、
$\vect{b}$ とその正射影のなす角 $\theta$ を使って $R^2 = \cos^2\theta$ と
書ける――第V部の「相関係数は $\cos\theta$」の続編である。
説明変数を追加することは部分空間 $W$ を広げることだから、
$R^2$ が決して下がらないことも幾何的に自明になる
(近い点の候補が増えるだけ)。
数式の海に見える計量経済学の背後に、つねに
「正射影」という一枚の絵があることを覚えておくと、見通しが劇的によくなる。
\end{itembox}

\begin{mondai}\label{q:lsq}
データ $(0, 0), (1, 1), (2, 1)$ に直線 $y = a + bx$ を
最小二乗法で当てはめよ。また、残差ベクトルが
$\vect{1} = \mat{1 \\ 1 \\ 1}$ と直交することを確かめよ。
\end{mondai}

\section{当てはめは直線に限らない}

「直線の当てはめ」と聞くと適用範囲が狭く感じるかもしれないが、
最小二乗法が要求するのは、モデルが\textbf{未知パラメータについて}
線形であることだけである。たとえば放物線 $y = a + bx + cx^2$ を
データに当てはめる問題は、各データ点 $(x_i, y_i)$ を代入した式
\[
a + b x_i + c x_i^2 = y_i \qquad (i = 1, \dots, m)
\]
を並べれば、未知数 $(a, b, c)$ についての連立\textbf{一次}方程式
$A\vect{x} = \vect{b}$ になる($A$ の第 $i$ 行は $(1,\ x_i,\ x_i^2)$)。
$x^2$ という非線形な項が入っていても、\textbf{係数 $a, b, c$ の入り方は
一次}だから、正規方程式がそのまま使える(例題 \ref{rei:lsq-parabola})。
$\sin x$ や $e^x$ を部品にした当てはめも同様である。
「変数について非線形でも、パラメータについて線形なら最小二乗法」
――適用範囲は見かけよりずっと広い。

\section{正射影行列 ― 「当てはめ」を一つの行列にまとめる}

列が線形独立な場合、最小二乗解は
$\vect{x}^* = ({}^t\!AA)^{-1}{}^t\!A\vect{b}$ と閉じた式で書けた。
このとき「最良の近似点」(回帰でいえば予測値のベクトル)は
\[
A\vect{x}^* = A({}^t\!AA)^{-1}{}^t\!A\,\vect{b}
\]
である。$\vect{b}$ の左に掛かっている行列に名前をつけよう。

\begin{teiri}[正射影行列]\label{thm:lsq-projmat}
$A$ を列が線形独立な $m \times n$ 行列とし、
\[
P = A({}^t\!AA)^{-1}{}^t\!A
\]
($m$ 次正方行列)とおく。このとき:
\begin{enumerate}
\item 任意の $\vect{b} \in \R^m$ に対し、$P\vect{b}$ は $\vect{b}$ の
$W = \mathrm{Im}\,A$ への正射影 $\vect{b}_W$ である。
\item ${}^t\!P = P$(対称)かつ $P^2 = P$(二度射影しても変わらない)。
\end{enumerate}
$P$ を $W$ への\textbf{正射影行列}という。
\end{teiri}

\noindent\textbf{(証明)}\quad
(1) $\vect{x}^* = ({}^t\!AA)^{-1}{}^t\!A\vect{b}$ は正規方程式の解だから、
定理 \ref{thm:lsq-normal} の証明で見たとおり
$A\vect{x}^* = \vect{b}_W$。左辺は $P\vect{b}$ そのものである。

(2) 転置の規則 ${}^t\!(XY) = {}^t Y\,{}^t\!X$ と
${}^t\!\bigl(({}^t\!AA)^{-1}\bigr) = \bigl({}^t\!({}^t\!AA)\bigr)^{-1}
= ({}^t\!AA)^{-1}$(対称行列の逆行列は対称)より
\[
{}^t\!P = {}^t\!\bigl(A({}^t\!AA)^{-1}{}^t\!A\bigr)
= A\,{}^t\!\bigl(({}^t\!AA)^{-1}\bigr)\,{}^t\!A = A({}^t\!AA)^{-1}{}^t\!A = P.
\]
また、内側の ${}^t\!A A$ が $({}^t\!AA)^{-1}$ と打ち消し合って
\[
P^2 = A({}^t\!AA)^{-1}\underbrace{{}^t\!A\,A({}^t\!AA)^{-1}}_{=\,E}{}^t\!A
= A({}^t\!AA)^{-1}{}^t\!A = P. \qquad\blacksquare
\]

$P^2 = P$ は「一度 $W$ に落とした点は、もう一度落としても動かない」
という正射影の幾何を式にしたものである(第V部のスペクトル分解に現れた
正射影行列 $\vect{u}\,{}^t\vect{u}$ は、$n = 1$ 本の場合の $P$ に
ほかならない)。統計学ではこの $P$ を\textbf{ハット行列}と呼ぶ
――データ $\vect{b} = (y_1, \dots, y_m)$ に掛けると予測値
$\hat{y}_1, \dots, \hat{y}_m$(ハットつきの $y$)が出てくるからである。

\section{反例 ― ありがちな誤解を正す}

\begin{hanrei}[列が線形従属なら最小二乗解は一つに定まらない]\label{hanrei:lsq-dependent}
「正規方程式は必ず解ける」と「解がただ一つ」は別の主張である。
\[
A = \mat{1 & 1 \\ 1 & 1}, \qquad \vect{b} = \mat{1 \\ 0}
\]
では、$A$ の列が線形従属で、
\[
{}^t\!AA = \mat{2 & 2 \\ 2 & 2}, \qquad {}^t\!A\vect{b} = \mat{1 \\ 1}
\]
より正規方程式は $2x + 2y = 1$(2本とも同じ式)。よって最小二乗解は
直線 $x + y = \dfrac{1}{2}$ 上の\textbf{全点}であり、無数にある。
実際 $A\vect{x} = (x + y)\mat{1 \\ 1}$ だから、
$\vect{b}$ との距離は和 $x + y$ だけで決まり、個々の $x, y$ には
依存しない。${}^t\!AA$ が正則になる(解が一意に定まる)のは、
定理 \ref{thm:lsq-normal} のとおり\textbf{列が線形独立のとき}に限る。
\end{hanrei}

\begin{hanrei}[「最良の直線」は当てはめの向きに依存する]\label{hanrei:lsq-swap}
最小二乗法の回帰直線を「点たちに最も寄り添う、ただ一本の直線」と
思ってはいけない。最小化しているのはあくまで\textbf{$y$ 方向の}残差
(縦のずれ)である。データ $(0, 0), (1, 0), (1, 1)$ で確かめよう。
$y = a + bx$ を当てはめると
\[
{}^t\!AA = \mat{3 & 2 \\ 2 & 2}, \quad {}^t\!A\vect{b} = \mat{1 \\ 1}
\quad\Longrightarrow\quad a = 0,\ b = \frac{1}{2}
\]
で $y = \dfrac{1}{2}x$。今度は役割を入れ替えて $x = c + dy$
(横のずれを最小化)とすると、データは $(y, x)$ の組として
$(0,0), (0,1), (1,1)$ だから
\[
\mat{3 & 1 \\ 1 & 1}\mat{c \\ d} = \mat{2 \\ 1}
\quad\Longrightarrow\quad c = d = \frac{1}{2}
\]
で $x = \dfrac{1}{2} + \dfrac{1}{2}y$、すなわち $y = 2x - 1$。
同じ $3$ 点から、\textbf{傾きが $4$ 倍違う二本の「最良の直線」}が
出てきた。どちらも正しい――最小化する残差の向きが違うだけである。
「何を誤差とみなすか」を決めるのは数学ではなく、データの出どころ
(どちらの変数に測定誤差があるか)である。
\end{hanrei}

\section{例題}

正規方程式による当てはめと、正射影の幾何の使い方を練習する。
以下は\textbf{解答つきの例題}である。手を動かして追ってほしい。

\begin{rei}[例題・基本(平均は最小二乗解)]\label{rei:lsq-mean}
一つの量 $x$ を $3$ 回測定して $1, 2, 4$ を得た。
「$x = 1$, $x = 2$, $x = 4$」という解のない連立方程式の
最小二乗解を求めよ。\\
\textbf{解.}\ 行列で書くと
$A = \mat{1 \\ 1 \\ 1}$, $\vect{b} = \mat{1 \\ 2 \\ 4}$。
${}^t\!AA = 3$($1 \times 1$ 行列), ${}^t\!A\vect{b} = 1 + 2 + 4 = 7$
だから、正規方程式は $3x = 7$ より
\[
x^* = \frac{7}{3} = \frac{1 + 2 + 4}{3}.
\]
――\textbf{測定値の平均}である。「くり返し測定の平均をとる」という
日常の作法は、二乗誤差最小の意味での最良の推定だったのである。
残差は $\mat{1 - 7/3 \\ 2 - 7/3 \\ 4 - 7/3} = \dfrac{1}{3}\mat{-4 \\ -1 \\ 5}$
で、成分の和が $0$(列 $\vect{1}$ と直交)になっている。
\end{rei}

\begin{rei}[例題・基本(回帰直線)]\label{rei:lsq-line}
データ $(0, 1), (1, 3), (2, 4)$ に直線 $y = a + bx$ を当てはめよ。\\
\textbf{解.}\
\[
A = \mat{1 & 0 \\ 1 & 1 \\ 1 & 2}, \quad \vect{b} = \mat{1 \\ 3 \\ 4}, \quad
{}^t\!AA = \mat{3 & 3 \\ 3 & 5}, \quad {}^t\!A\vect{b} = \mat{8 \\ 11}.
\]
正規方程式 $3a + 3b = 8$, $3a + 5b = 11$ を引き算して $2b = 3$、
$b = \dfrac{3}{2}$, $a = \dfrac{7}{6}$。回帰直線は
$y = \dfrac{7}{6} + \dfrac{3}{2}x$ である。
検算:予測値は $\dfrac{7}{6}, \dfrac{8}{3}, \dfrac{25}{6}$、残差は
$\mat{-1/6 \\ 1/3 \\ -1/6}$ で、成分の和は $0$、
$x$ 座標との内積も $0 \cdot \left(-\dfrac{1}{6}\right)
+ 1 \cdot \dfrac{1}{3} + 2 \cdot \left(-\dfrac{1}{6}\right) = 0$。
残差が $A$ の両方の列と直交している。
\end{rei}

\begin{rei}[例題・基本(残差の直交性)]\label{rei:lsq-resid}
$A = \mat{1 & 0 \\ 0 & 1 \\ 1 & 1}$, $\vect{b} = \mat{1 \\ 1 \\ 1}$ の
最小二乗解を求め、残差が $A$ の各列と直交することを確かめよ。\\
\textbf{解.}\
${}^t\!AA = \mat{2 & 1 \\ 1 & 2}$, ${}^t\!A\vect{b} = \mat{2 \\ 2}$。
正規方程式を解いて $\vect{x}^* = \mat{2/3 \\ 2/3}$。このとき
\[
A\vect{x}^* = \mat{2/3 \\ 2/3 \\ 4/3}, \qquad
\vect{r} = \vect{b} - A\vect{x}^* = \frac{1}{3}\mat{1 \\ 1 \\ -1}.
\]
第 $1$ 列との内積は $\dfrac{1}{3}(1 + 0 - 1) = 0$、
第 $2$ 列との内積は $\dfrac{1}{3}(0 + 1 - 1) = 0$。
確かに残差は $\mathrm{Im}\,A$ と直交している
(これが「これ以上改善できない」ことの幾何的な証明書である)。
\end{rei}

\begin{rei}[例題・標準(放物線の当てはめ)]\label{rei:lsq-parabola}
データ $(-1, 2), (0, 1), (1, 4)$ に、モデル $y = a + bx^2$ を
最小二乗法で当てはめよ。\\
\textbf{解.}\ 各点を代入した式 $a + b x_i^2 = y_i$ を並べると
\[
A = \mat{1 & 1 \\ 1 & 0 \\ 1 & 1}, \qquad \vect{b} = \mat{2 \\ 1 \\ 4}
\]
($A$ の第 $2$ 列は $x_i^2$ の値)。モデルは $x$ については
$2$ 次だが、\textbf{未知数 $a, b$ については一次}なので正規方程式が使える:
\[
{}^t\!AA = \mat{3 & 2 \\ 2 & 2}, \qquad {}^t\!A\vect{b} = \mat{7 \\ 6}.
\]
$3a + 2b = 7$, $2a + 2b = 6$ を引き算して $a = 1$、$b = 2$。
当てはめた曲線は $y = 1 + 2x^2$ である。
検算:予測値は $3, 1, 3$、残差は $(-1, 0, 1)$ で、
$A$ の第 $1$ 列 $(1,1,1)$ とも第 $2$ 列 $(1,0,1)$ とも内積 $0$。
\end{rei}

\begin{rei}[例題・標準(正射影行列)]\label{rei:lsq-projmat-ex}
$A = \mat{1 \\ 1 \\ 1}$ に対する正射影行列 $P$ を求め、
$P\vect{b}$($\vect{b} = {}^t(1, 2, 4)$)を計算せよ。\\
\textbf{解.}\ ${}^t\!AA = 3$ より
\[
P = A \cdot \frac{1}{3} \cdot {}^t\!A
= \frac{1}{3}\mat{1 & 1 & 1 \\ 1 & 1 & 1 \\ 1 & 1 & 1}.
\]
${}^t\!P = P$ は見た目から明らか、$P^2 = P$ も
(各成分を計算すると $\frac{1}{9} \times 3 = \frac{1}{3}$)確かめられる。
$\vect{b} = {}^t(1, 2, 4)$ に施すと
\[
P\vect{b} = \frac{1}{3}\mat{7 \\ 7 \\ 7}
\]
――すべての成分が\textbf{平均 $7/3$} のベクトルになる。
例題 \ref{rei:lsq-mean} の「平均は最小二乗解」を、
「$\vect{b}$ を対角線方向 $\vect{1}$ の直線に正射影すると
平均のベクトルになる」と言い換えたことになる。
\end{rei}

\begin{rei}[例題・標準(証明:${}^t\!AA$ と $A$ の核は一致する)]\label{rei:lsq-kernel}
任意の $m \times n$ 行列 $A$ について
$\mathrm{Ker}({}^t\!AA) = \mathrm{Ker}\,A$ を示し、
$\mathrm{rank}({}^t\!AA) = \mathrm{rank}\,A$ を導け。\\
\textbf{解.}\ $A\vect{x} = \vect{0}$ なら
${}^t\!AA\vect{x} = {}^t\!A\vect{0} = \vect{0}$ だから
$\mathrm{Ker}\,A \subseteq \mathrm{Ker}({}^t\!AA)$。
逆に ${}^t\!AA\vect{x} = \vect{0}$ なら、左から ${}^t\vect{x}$ を掛けて
\[
0 = {}^t\vect{x}\,{}^t\!AA\vect{x} = \|A\vect{x}\|^2
\]
より $A\vect{x} = \vect{0}$(ノルムの正値性)。よって両者は一致する。
${}^t\!AA$ も $A$ も列の本数は $n$ だから、次元定理
(定理 \ref{thm:dimension})より
\[
\mathrm{rank}({}^t\!AA) = n - \dim\mathrm{Ker}({}^t\!AA)
= n - \dim\mathrm{Ker}\,A = \mathrm{rank}\,A.
\]
「${}^t\!A$ を掛けても階数の情報は失われない」――この事実は
次章以降(QR分解・特異値分解)でもくり返し効いてくる。
\end{rei}

\begin{matome}
\begin{itemize}
\item \textbf{最小二乗問題}:$A\vect{x} = \vect{b}$ が解けないとき、
残差のノルム $\|\vect{b} - A\vect{x}\|$ を最小にする $\vect{x}$ を探す。
答えは「$\vect{b}$ を $W = \mathrm{Im}\,A$ に正射影せよ」
(第V部の最良近似)。
\item 最小二乗解 $\iff$ \textbf{正規方程式} ${}^t\!AA\vect{x} = {}^t\!A\vect{b}$
(定理 \ref{thm:lsq-normal})。正規方程式は必ず解をもち、
$A$ の列が線形独立なら解はただ一つ
$\vect{x}^* = ({}^t\!AA)^{-1}{}^t\!A\vect{b}$。
\item 残差 $\vect{b} - A\vect{x}^*$ は $\mathrm{Im}\,A$ と\textbf{直交}する
(検算にも使える)。回帰直線・多項式当てはめなど、
\textbf{パラメータについて線形}なモデルはすべてこの枠組みで解ける。
\item \textbf{正射影行列} $P = A({}^t\!AA)^{-1}{}^t\!A$ は
$\vect{b} \mapsto \vect{b}_W$ を与え、${}^t\!P = P$, $P^2 = P$ を満たす
(定理 \ref{thm:lsq-projmat})。
\item 列が線形従属なら最小二乗解は無数にある(反例 \ref{hanrei:lsq-dependent})。
また回帰直線は $y$ 方向の残差を最小化するもので、$x$ と $y$ の
役割を入れ替えると別の直線になる(反例 \ref{hanrei:lsq-swap})。
\end{itemize}
\end{matome}

\section*{章末問題}

\begin{mondai}[基本]\label{q:lsq-ex1}
一つの量を $3$ 回測定して $2, 3, 7$ を得た。解のない連立方程式
「$x = 2$, $x = 3$, $x = 7$」の最小二乗解を正規方程式で求め、
残差ベクトルの成分の和が $0$ になることを確かめよ。
\end{mondai}

\begin{mondai}[基本]\label{q:lsq-ex2}
データ $(0, 0), (1, 2), (2, 3)$ に直線 $y = a + bx$ を
最小二乗法で当てはめよ。
\end{mondai}

\begin{mondai}[基本]\label{q:lsq-ex3}
連立方程式
\[
x + y = 1, \qquad x - y = 1, \qquad x + y = 3
\]
は解をもたない。最小二乗解を求め、残差ベクトルが係数行列の
$2$ つの列と直交することを確かめよ。
\end{mondai}

\begin{mondai}[基本]\label{q:lsq-ex4}
$A = \mat{1 & 0 \\ 1 & 1 \\ 1 & 2}$, $\vect{b} = \mat{1 \\ 2 \\ 3}$ について
正規方程式を解き、得られた $\vect{x}^*$ が $A\vect{x} = \vect{b}$ の
\textbf{厳密な}解になっている(残差が $\vect{0}$)ことを確かめよ。
データの言葉でいえば、$3$ 点 $(0,1), (1,2), (2,3)$ が
一直線上に並んでいる場合である。
\end{mondai}

\begin{mondai}[標準]\label{q:lsq-ex5}
データ $(0, 0), (1, 1), (2, 3), (3, 4)$ に直線 $y = a + bx$ を
最小二乗法で当てはめ、残差ベクトルを求めて検算せよ。
\end{mondai}

\begin{mondai}[標準]\label{q:lsq-ex6}
$A = \mat{1 & 0 \\ 0 & 1 \\ 1 & 1}$(例題 \ref{rei:lsq-resid} の行列)への
正射影行列 $P = A({}^t\!AA)^{-1}{}^t\!A$ を求め、
${}^t\!P = P$ と $P^2 = P$ を確かめよ。さらに
$P\mat{1 \\ 1 \\ 1}$ が例題 \ref{rei:lsq-resid} の $A\vect{x}^*$ と
一致することを確かめよ。
\end{mondai}

\begin{mondai}[標準]\label{q:lsq-ex7}
問 \ref{q:lsq} のデータ $(0, 0), (1, 1), (2, 1)$ について、今度は
$x$ と $y$ の役割を入れ替え、$x = c + dy$ を最小二乗法で当てはめよ。
得られる直線を $y$ について解いた形に直し、問 \ref{q:lsq} の
回帰直線 $y = \dfrac{1}{6} + \dfrac{1}{2}x$ と比較せよ。
\end{mondai}

\begin{mondai}[標準]\label{q:lsq-ex8}
原点を通る直線 $y = kx$ をデータ $(1, 2), (2, 3), (3, 4)$ に
最小二乗法で当てはめよ(未知数は $k$ の一つだけである)。
また、残差ベクトルが $\mat{1 \\ 2 \\ 3}$ と直交することを確かめよ。
\end{mondai}

\begin{mondai}[発展]\label{q:lsq-ex9}
$A = \mat{1 & 2 \\ 2 & 4 \\ 0 & 0}$, $\vect{b} = \mat{1 \\ 1 \\ 1}$ について、\\
　(1) 正規方程式を書き下し、最小二乗解の全体が直線
$x + 2y = \dfrac{3}{5}$ になることを示せ。\\
　(2) その中でノルム $\|(x, y)\|$ が最小のものを求めよ
(この「ノルム最小の最小二乗解」が第 \ref{chap:pinv} 章の主役になる)。
\end{mondai}

\begin{mondai}[発展]\label{q:lsq-ex10}
データ $(x_1, y_1), \dots, (x_m, y_m)$ への直線 $y = a + bx$ の
当てはめを一般に考える。
$S_x = \sum x_i$, $S_y = \sum y_i$, $S_{xx} = \sum x_i^2$,
$S_{xy} = \sum x_i y_i$ とおくとき、正規方程式が
\[
\mat{m & S_x \\ S_x & S_{xx}}\mat{a \\ b} = \mat{S_y \\ S_{xy}}
\]
となることを示し、$m S_{xx} - S_x^2 \neq 0$ のとき
\[
b = \frac{m S_{xy} - S_x S_y}{m S_{xx} - S_x^2}, \qquad
a = \bar{y} - b\bar{x}
\quad \left(\bar{x} = \frac{S_x}{m},\ \bar{y} = \frac{S_y}{m}\right)
\]
を導け(統計学の教科書に載っている回帰係数の公式である)。
問 \ref{q:lsq} のデータで検算せよ。
\end{mondai}

\begin{mondai}[証明]\label{q:lsq-ex11}
任意の $m \times n$ 行列 $A$ について、${}^t\!AA$ は対称かつ半正定値
(第V部)であることを示せ。さらに、${}^t\!AA$ が正定値であることと
$A$ の列が線形独立であることが同値であることを示せ。
\end{mondai}

\begin{mondai}[証明]\label{q:lsq-ex12}
$W = \mathrm{Im}\,A$ とし、$\vect{b} = \vect{b}_W + \vect{b}_{\perp}$ を
直交分解(定理 \ref{thm:gs-decomp})とする。任意の $\vect{x}$ に対して
\[
\|\vect{b} - A\vect{x}\|^2 = \|\vect{b}_{\perp}\|^2 + \|\vect{b}_W - A\vect{x}\|^2
\]
が成り立つことを示し、最小二乗解の全体が
「$A\vect{x} = \vect{b}_W$ の解の全体」と一致することを結論せよ。
\end{mondai}

\bigskip
{\small
\noindent\textbf{【章末問題の方針】}\quad
まず自力で取り組み、行き詰まったら下の方針を見よ。記述による完全解答は
巻末「問の解答」にある。
\begin{itemize}
\item \textbf{問 \ref{q:lsq-ex1}}:$A = {}^t(1, 1, 1)$。正規方程式は
$3x = (\text{測定値の和})$。
\item \textbf{問 \ref{q:lsq-ex2}}:${}^t\!AA$, ${}^t\!A\vect{b}$ を作って
$2$ 元連立を解く。例題 \ref{rei:lsq-line} と同じ手順。
\item \textbf{問 \ref{q:lsq-ex3}}:$A = \mat{1 & 1 \\ 1 & -1 \\ 1 & 1}$,
$\vect{b} = {}^t(1, 1, 3)$ とおいて正規方程式。
\item \textbf{問 \ref{q:lsq-ex4}}:正規方程式を解くと残差が $\vect{0}$ になる。
「解があるときは最小二乗解 $=$ 通常の解」の確認である。
\item \textbf{問 \ref{q:lsq-ex5}}:${}^t\!AA = \mat{4 & 6 \\ 6 & 14}$。
残差の成分の和と、$x$ 座標との内積がともに $0$ なら検算完了。
\item \textbf{問 \ref{q:lsq-ex6}}:$({}^t\!AA)^{-1} = \dfrac{1}{3}
\mat{2 & -1 \\ -1 & 2}$。$P$ は $3$ 次正方行列になる。
\item \textbf{問 \ref{q:lsq-ex7}}:データを $(y, x)$ の組と読み替えて
同じ手順を踏む。傾きの逆数にはならないことに注意。
\item \textbf{問 \ref{q:lsq-ex8}}:$A = {}^t(1, 2, 3)$ で ${}^t\!AA = 14$。
正規方程式は $1$ 本である。
\item \textbf{問 \ref{q:lsq-ex9}}:列が従属なので反例
\ref{hanrei:lsq-dependent} と同じ状況。(2) は原点から直線への垂線の足。
\item \textbf{問 \ref{q:lsq-ex10}}:$A$ の第 $i$ 行は $(1, x_i)$。
${}^t\!AA$ の成分を書き出すだけである。$b$ はクラメルの公式(第II部)か
消去法で。
\item \textbf{問 \ref{q:lsq-ex11}}:対称性は転置の規則。
${}^t\vect{x}({}^t\!AA)\vect{x} = \|A\vect{x}\|^2$ が鍵
(例題 \ref{rei:lsq-kernel} と同じ計算)。
\item \textbf{問 \ref{q:lsq-ex12}}:$\vect{b} - A\vect{x} =
\vect{b}_{\perp} + (\vect{b}_W - A\vect{x})$ と分けてピタゴラス
(定理 \ref{thm:inner-pyth})。等号成立の条件を読む。
\end{itemize}
}

% =====================================================
\chapter{QR分解 ― グラム・シュミットの行列表現}\label{chap:qr}
% =====================================================

\begin{goal}
\begin{itemize}
\item 列が線形独立な行列が $A = QR$($Q$ の列は正規直交系、$R$ は
対角成分が正の上三角行列)と\textbf{ただ一通り}に分解できることを、
存在(グラム・シュミットの計算記録)と一意性の証明つきで理解する。
\item $2 \times 2$ や $3 \times 2$ の行列の QR 分解を手計算で
求められるようになる。
\item QR 分解を使って最小二乗問題を $R\vect{x} = {}^t\!Q\vect{b}$ の
\textbf{後退代入}に帰着できるようになる。
\item 正方行列では $Q$ が直交行列になり、
$|\det A| = r_{11}r_{22}\cdots r_{nn}$ が成り立つことを使えるようになる。
\end{itemize}
\end{goal}

\section{動機 ― 「計算の記録」を行列の等式にする}

第V部のグラム・シュミットの直交化法は、与えられた列ベクトルの組を
正規直交系に作り替える\textbf{手続き}だった。本章はその手続きの
計算結果を、$A = QR$ という一本の\textbf{行列の等式}として書き留める。
たったそれだけのことが、驚くほど役に立つ。

本書ではすでに、行列を「よい形の積」に割る操作に何度も出会っている。
対角化 $A = PDP^{-1}$(第IV部)、直交対角化 $A = U D\,{}^t U$(第V部)、
ジョルダン標準形(第VI部)――いずれも、複雑な行列を
「性質のよくわかる部品の積」に分解することで、べき乗や微分方程式が
一気に解けた。このような書き換えを\textbf{行列の分解}と総称する。
QR分解はその仲間で、部品は「直交(回転・鏡映の仲間)」と
「上三角(下から順に解ける)」である。理論の主役というより
\textbf{計算の主役}であり、次章の特異値分解と並んで、
数値計算ソフトウェアの心臓部で毎日動いている。

\section{QR分解 ― 存在と一意性}

\begin{teiri}[QR分解]\label{thm:qr-main}
列が線形独立な $m \times n$ 行列 $A$ は、
\[
A = QR
\]
($Q$ は列が正規直交系の $m \times n$ 行列、$R$ は対角成分が正の
$n$ 次上三角行列)の形に\textbf{ただ一通り}に分解できる。
\end{teiri}

新しい定理に見えるが、正体は\textbf{グラム・シュミットの直交化法の
計算記録}である。$A$ の列を直交化して得た正規直交系を $Q$ の列とし、
「元の各列が、直交化済みベクトルのどんな線形結合だったか」の係数を
$R$ に書き留める。$j$ 番目の列は $1, \dots, j$ 番目の直交ベクトルだけで
書けるから、$R$ は自動的に上三角になる。以下、これを証明の形に直す。

\noindent\textbf{(証明)}\quad
\textbf{(存在)}\ $A$ の列を $\vect{a}_1, \dots, \vect{a}_n$ とし、
グラム・シュミットの直交化法(定理 \ref{thm:gs-main})を適用して
正規直交系 $\vect{u}_1, \dots, \vect{u}_n$ を作る。手続きの式
\[
\vect{b}_j = \vect{a}_j - \sum_{i=1}^{j-1}
\langle \vect{a}_j, \vect{u}_i \rangle \vect{u}_i, \qquad
\vect{u}_j = \frac{\vect{b}_j}{\|\vect{b}_j\|}
\]
を $\vect{a}_j$ について解き直すと
\[
\vect{a}_j = \langle \vect{a}_j, \vect{u}_1 \rangle \vect{u}_1 + \cdots
+ \langle \vect{a}_j, \vect{u}_{j-1} \rangle \vect{u}_{j-1}
+ \|\vect{b}_j\|\,\vect{u}_j.
\]
そこで $Q = (\vect{u}_1 \ \cdots \ \vect{u}_n)$、$R$ の $(i, j)$ 成分を
\[
r_{ij} = \langle \vect{a}_j, \vect{u}_i \rangle \ (i < j), \qquad
r_{jj} = \|\vect{b}_j\| \ (> 0), \qquad
r_{ij} = 0 \ (i > j)
\]
と定めれば、上の式は「$A$ の第 $j$ 列 $=$ $QR$ の第 $j$ 列」を
意味するから $A = QR$ である($\vect{b}_j \neq \vect{0}$、すなわち
$r_{jj} > 0$ が保証されるのは列が線形独立だからである。
反例 \ref{hanrei:gs-dependent} 参照)。

\textbf{(一意性)}\ $A = QR = Q'R'$ と二通りに分解できたとし、
$Q, Q'$ の列を $\vect{q}_i, \vect{q}'_i$、$R, R'$ の成分を
$r_{ij}, r'_{ij}$ とする。左の列から順に一致を示す。
第 $1$ 列を比べると $\vect{a}_1 = r_{11}\vect{q}_1 = r'_{11}\vect{q}'_1$。
ノルムをとると($\|\vect{q}_1\| = \|\vect{q}'_1\| = 1$ より)
$r_{11} = r'_{11} = \|\vect{a}_1\|$、よって $\vect{q}_1 = \vect{q}'_1$。
次に、第 $j - 1$ 列までで $\vect{q}_i = \vect{q}'_i$($i < j$)が
示せたとする。第 $j$ 列は
\[
\vect{a}_j = \sum_{i=1}^{j-1} r_{ij}\vect{q}_i + r_{jj}\vect{q}_j
= \sum_{i=1}^{j-1} r'_{ij}\vect{q}_i + r'_{jj}\vect{q}'_j.
\]
$\vect{q}_i$($i < j$)との内積をとると、$\vect{q}_j$ も $\vect{q}'_j$ も
$\vect{q}_i$ と直交する($Q'$ の列どうしの直交性と
$\vect{q}_i = \vect{q}'_i$)から、$r_{ij} = \langle \vect{a}_j, \vect{q}_i
\rangle = r'_{ij}$。残った項を比べて
$r_{jj}\vect{q}_j = r'_{jj}\vect{q}'_j$、ノルムをとって
$r_{jj} = r'_{jj} > 0$、よって $\vect{q}_j = \vect{q}'_j$。
帰納的にすべての列が一致し、$Q = Q'$, $R = R'$ である。 $\blacksquare$

\begin{rei}[$2 \times 2$ の QR 分解]\label{rei:qr-2x2}
$A = \mat{3 & 1 \\ 4 & 3}$ を分解する。
$\vect{a}_1 = \mat{3 \\ 4}$ のノルムは $5$ だから
$\vect{u}_1 = \dfrac{1}{5}\mat{3 \\ 4}$, $r_{11} = 5$。
$r_{12} = \langle \vect{a}_2, \vect{u}_1 \rangle
= \dfrac{3 + 12}{5} = 3$ を使って
\[
\vect{b}_2 = \mat{1 \\ 3} - 3 \cdot \frac{1}{5}\mat{3 \\ 4}
= \frac{1}{5}\mat{-4 \\ 3}, \qquad r_{22} = \|\vect{b}_2\| = 1.
\]
よって
\[
A = QR = \frac{1}{5}\mat{3 & -4 \\ 4 & 3}\mat{5 & 3 \\ 0 & 1}.
\]
検算:$QR$ の第 $2$ 列は
$3\vect{u}_1 + 1 \cdot \vect{u}_2 = \dfrac{1}{5}\mat{9 - 4 \\ 12 + 3}
= \mat{1 \\ 3}$ で一致する。$Q$ は回転行列
(第V部:$\det Q = \frac{9 + 16}{25} = 1$)である。
\end{rei}

\section{QR分解で最小二乗法を解く}

QR分解を使うと、前章の最小二乗問題が見通しよく解ける。
$A = QR$ を正規方程式に代入すると ${}^t\!Q Q = E$ により
\[
{}^t\!R R \vect{x} = {}^t\!R\,{}^t\!Q\,\vect{b}
\quad\Longrightarrow\quad
R\vect{x} = {}^t\!Q\,\vect{b}
\]
となり($R$ は正則なので ${}^t\!R$ で割ってよい)、$R$ が上三角なので
\textbf{下から順に代入するだけ}で解ける。逆行列 $({}^t\!AA)^{-1}$ を
作る必要がそもそもなくなるのである。

\begin{rei}[QR 経由の最小二乗法]\label{rei:qr-lsq}
$A = \mat{1 & 1 \\ 1 & 0 \\ 0 & 1}$, $\vect{b} = \mat{1 \\ 1 \\ 1}$ の
最小二乗解を求める。直交化すると
\[
\vect{u}_1 = \frac{1}{\sqrt{2}}\mat{1 \\ 1 \\ 0}, \quad
\vect{b}_2 = \mat{1 \\ 0 \\ 1} - \frac{1}{2}\mat{1 \\ 1 \\ 0}
= \frac{1}{2}\mat{1 \\ -1 \\ 2}, \quad
\vect{u}_2 = \frac{1}{\sqrt{6}}\mat{1 \\ -1 \\ 2}
\]
で、$r_{11} = \sqrt{2}$, $r_{12} = \dfrac{1}{\sqrt{2}}$,
$r_{22} = \dfrac{\sqrt{6}}{2}$。右辺は
\[
{}^t\!Q\vect{b} = \mat{\langle \vect{b}, \vect{u}_1 \rangle \\
\langle \vect{b}, \vect{u}_2 \rangle}
= \mat{2/\sqrt{2} \\ 2/\sqrt{6}}.
\]
$R\vect{x} = {}^t\!Q\vect{b}$ を下から解くと、第 $2$ 行
$\dfrac{\sqrt{6}}{2}x_2 = \dfrac{2}{\sqrt{6}}$ より $x_2 = \dfrac{2}{3}$、
第 $1$ 行 $\sqrt{2}\,x_1 + \dfrac{1}{\sqrt{2}} \cdot \dfrac{2}{3}
= \dfrac{2}{\sqrt{2}}$ の両辺を $\sqrt{2}$ で割って
$x_1 + \dfrac{1}{3} = 1$、$x_1 = \dfrac{2}{3}$。
最小二乗解は $\vect{x}^* = {}^t\!\left(\dfrac{2}{3},\ \dfrac{2}{3}\right)$
(正規方程式 ${}^t\!AA = \mat{2 & 1 \\ 1 & 2}$,
${}^t\!A\vect{b} = \mat{2 \\ 2}$ を解いても同じ答えになる)。
\end{rei}

\begin{itembox}[l]{深掘り:数値計算の世界では QR が主役である}
理論上は正規方程式 ${}^t\!AA\vect{x} = {}^t\!A\vect{b}$ で十分だが、
コンピュータの有限桁計算では事情が変わる。${}^t\!AA$ を作る操作は
問題の「数値的な敏感さ」(条件数と呼ばれる量)を\textbf{二乗}してしまい、
桁落ちで精度が壊れることがある。QR分解はこの悪化を起こさないため、
実用ソフトウェア(統計パッケージの回帰計算など)の内部では
QR 経由の解法が標準である。さらに、QR分解を反復する
\textbf{QR法}は、固有値を数値的に求める現代の標準アルゴリズムでもある
(固有方程式を解く方法は $5$ 次以上で代数的に破綻するし、数値的にも悪い。
その入り口を章末問題 \ref{q:qr-ex10} で体験できる)。
第II部の深掘り「定義と計算法は別物」の好例が、ここにもある。
\end{itembox}

\begin{mondai}\label{q:qr}
$A = \mat{1 & 1 \\ 0 & 1 \\ 1 & 0}$ の QR 分解を求めよ
(列の直交化は問 \ref{q:gs} で実行済みである。それを流用してよい)。
\end{mondai}

\section{正方行列の場合 ― 直交行列 × 上三角、そして行列式}

$A$ が $n$ 次正方で正則(列が線形独立)なら、$Q$ も $n$ 次正方になり、
列が正規直交系の正方行列、すなわち\textbf{直交行列}である(第V部)。
「どんな正則行列も、直交行列と上三角行列の積」――これは
行列式とも美しく噛み合う。

\begin{teiri}[QR分解と行列式]\label{thm:qr-det}
正則な $n$ 次正方行列 $A = QR$ について
\[
|\det A| = r_{11}\,r_{22} \cdots r_{nn}.
\]
\end{teiri}

\noindent\textbf{(証明)}\quad
積の行列式(第II部)より $\det A = \det Q \cdot \det R$。
直交行列では ${}^t\!QQ = E$ の両辺の行列式をとると
$(\det Q)^2 = 1$、すなわち $\det Q = \pm 1$。また上三角行列の行列式は
対角成分の積 $r_{11}\cdots r_{nn}$(第II部)で、これは正である。
絶対値をとって主張を得る。 $\blacksquare$

第II部で「行列式は平行体の体積(の符号つき版)」と学んだ。
この定理はその言い換えである:$r_{jj} = \|\vect{b}_j\|$ は
「$j$ 番目の辺のうち、それまでの辺と直交する成分(高さ)」だから、
$\prod r_{jj}$ は「底辺 × 高さ × …」をくり返した体積計算に
ほかならない。行列式・直交化・体積という三つの話題が、
QR分解の上で一つに結ばれる(章末問題 \ref{q:qr-ex12} では、
ここから\textbf{アダマールの不等式}が一行で出ることを見る)。

\section{反例 ― ありがちな誤解を正す}

\begin{hanrei}[列が線形従属だと、この形の QR 分解は存在しない]\label{hanrei:qr-dependent}
$A = \mat{1 & 2 \\ 1 & 2}$(第 $2$ 列は第 $1$ 列の $2$ 倍)が
$A = QR$($Q$ の列は正規直交系、$R$ は対角成分が\textbf{正}の上三角)と
分解できたとする。$Q$ の列は正規直交系だから線形独立で、$R$ は
対角成分が正の三角行列だから正則。すると $A$ の列は
「線形独立なベクトルの組に正則行列を掛けたもの」となり
線形独立のはず――しかし実際は従属であり、矛盾する。
グラム・シュミットの手続きの側でも、$\vect{b}_2 = \vect{0}$ となって
正規化で止まる(反例 \ref{hanrei:gs-dependent})。
定理 \ref{thm:qr-main} の仮定「列が線形独立」は飾りではない。
\end{hanrei}

\begin{hanrei}[「対角成分が正」を外すと一意性が崩れる]\label{hanrei:qr-sign}
単位行列 $E$($2$ 次)は
\[
E = E \cdot E
= \mat{-1 & 0 \\ 0 & 1}\mat{-1 & 0 \\ 0 & 1}
\]
と\textbf{二通り}に「直交行列 × 上三角行列」に書ける
($\mat{-1 & 0 \\ 0 & 1}$ は鏡映で直交行列、同時に対角行列なので
上三角でもある)。二つ目の分解は $r_{11} = -1 < 0$ なので、
定理 \ref{thm:qr-main} の条件「対角成分が正」に反する。
つまりこの条件は、各 $\vect{u}_j$ の\textbf{向きの選び方}
($\pm$ の任意性)を一つに固定するための約束であり、
これを外すと分解は($\pm$ の分だけ)複数現れる。
一意性を支えているのが不等式 $r_{jj} > 0$ だという点は、
証明の「ノルムをとって $r_{jj} = r'_{jj}$」の段で確かめられる。
\end{hanrei}

\section{例題}

QR 分解の計算と、最小二乗法・行列式への応用を練習する。
以下は\textbf{解答つきの例題}である。手を動かして追ってほしい。

\begin{rei}[例題・基本($2 \times 2$ の分解)]\label{rei:qr-basic2x2}
$A = \mat{1 & 2 \\ 1 & 0}$ の QR 分解を求めよ。\\
\textbf{解.}\ $\vect{a}_1 = \mat{1 \\ 1}$ より $r_{11} = \sqrt{2}$,
$\vect{u}_1 = \dfrac{1}{\sqrt{2}}\mat{1 \\ 1}$。
$r_{12} = \langle \vect{a}_2, \vect{u}_1 \rangle = \dfrac{2}{\sqrt{2}}
= \sqrt{2}$、
\[
\vect{b}_2 = \mat{2 \\ 0} - \sqrt{2} \cdot \frac{1}{\sqrt{2}}\mat{1 \\ 1}
= \mat{1 \\ -1}, \qquad r_{22} = \sqrt{2},\quad
\vect{u}_2 = \frac{1}{\sqrt{2}}\mat{1 \\ -1}.
\]
よって
\[
A = \frac{1}{\sqrt{2}}\mat{1 & 1 \\ 1 & -1}\mat{\sqrt{2} & \sqrt{2} \\ 0 & \sqrt{2}}.
\]
検算:第 $2$ 列は $\sqrt{2}\vect{u}_1 + \sqrt{2}\vect{u}_2
= \mat{1 \\ 1} + \mat{1 \\ -1} = \mat{2 \\ 0}$。
また $|\det A| = |0 - 2| = 2 = \sqrt{2} \cdot \sqrt{2}$
(定理 \ref{thm:qr-det})。
\end{rei}

\begin{rei}[例題・基本($3 \times 2$ の分解)]\label{rei:qr-basic3x2}
$A = \mat{1 & 2 \\ 0 & 3 \\ 0 & 4}$ の QR 分解を求めよ。\\
\textbf{解.}\ $\vect{a}_1 = \vect{e}_1$ はすでに単位ベクトルなので
$r_{11} = 1$, $\vect{u}_1 = \vect{e}_1$。
$r_{12} = \langle \vect{a}_2, \vect{u}_1 \rangle = 2$、
\[
\vect{b}_2 = \mat{2 \\ 3 \\ 4} - 2\mat{1 \\ 0 \\ 0} = \mat{0 \\ 3 \\ 4},
\qquad r_{22} = 5, \quad \vect{u}_2 = \frac{1}{5}\mat{0 \\ 3 \\ 4}.
\]
よって
\[
A = QR = \mat{1 & 0 \\ 0 & 3/5 \\ 0 & 4/5}\mat{1 & 2 \\ 0 & 5}.
\]
「第 $1$ 列の方向を残し、第 $2$ 列からその成分を取り除いて正規化した」
という直交化の履歴が、$R$ にそのまま記録されている。
\end{rei}

\begin{rei}[例題・基本(QR で最小二乗)]\label{rei:qr-lsq-ex}
例題 \ref{rei:qr-basic3x2} の $A$ と $\vect{b} = \mat{1 \\ 5 \\ 5}$ に
対する最小二乗解を、$R\vect{x} = {}^t\!Q\vect{b}$ で求めよ。\\
\textbf{解.}\
\[
{}^t\!Q\vect{b} = \mat{\langle \vect{b}, \vect{u}_1 \rangle \\
\langle \vect{b}, \vect{u}_2 \rangle}
= \mat{1 \\ (15 + 20)/5} = \mat{1 \\ 7}.
\]
$R\vect{x} = \mat{1 \\ 7}$ を下から解く:第 $2$ 行 $5x_2 = 7$ より
$x_2 = \dfrac{7}{5}$、第 $1$ 行 $x_1 + 2x_2 = 1$ より
$x_1 = 1 - \dfrac{14}{5} = -\dfrac{9}{5}$。
検算:正規方程式は ${}^t\!AA = \mat{1 & 2 \\ 2 & 29}$,
${}^t\!A\vect{b} = \mat{1 \\ 37}$ で、
$2 \cdot \left(-\dfrac{9}{5}\right) + 29 \cdot \dfrac{7}{5}
= \dfrac{-18 + 203}{5} = 37$ が成り立つ。
\end{rei}

\begin{rei}[例題・標準($3 \times 3$ の分解と行列式)]\label{rei:qr-3x3}
$A = \mat{1 & 1 & 0 \\ 1 & 0 & 1 \\ 0 & 1 & 1}$ の QR 分解を求め、
定理 \ref{thm:qr-det} を検証せよ。\\
\textbf{解.}\ 第 $1$・第 $2$ 列の直交化は例題 \ref{rei:qr-lsq} と同じ計算で
\[
\vect{u}_1 = \frac{1}{\sqrt{2}}\mat{1 \\ 1 \\ 0}, \qquad
\vect{u}_2 = \frac{1}{\sqrt{6}}\mat{1 \\ -1 \\ 2}, \qquad
r_{11} = \sqrt{2},\ r_{12} = \frac{1}{\sqrt{2}},\ r_{22} = \frac{\sqrt{6}}{2}.
\]
第 $3$ 列 $\vect{a}_3 = {}^t(0, 1, 1)$ は
$r_{13} = \langle \vect{a}_3, \vect{u}_1 \rangle = \dfrac{1}{\sqrt{2}}$,
$r_{23} = \langle \vect{a}_3, \vect{u}_2 \rangle = \dfrac{1}{\sqrt{6}}$ で、
\[
\vect{b}_3 = \mat{0 \\ 1 \\ 1} - \frac{1}{2}\mat{1 \\ 1 \\ 0}
- \frac{1}{6}\mat{1 \\ -1 \\ 2} = \frac{2}{3}\mat{-1 \\ 1 \\ 1},
\qquad r_{33} = \|\vect{b}_3\| = \frac{2}{\sqrt{3}}, \quad
\vect{u}_3 = \frac{1}{\sqrt{3}}\mat{-1 \\ 1 \\ 1}.
\]
よって $Q = (\vect{u}_1\ \vect{u}_2\ \vect{u}_3)$,
$R = \mat{\sqrt{2} & 1/\sqrt{2} & 1/\sqrt{2} \\
0 & \sqrt{6}/2 & 1/\sqrt{6} \\ 0 & 0 & 2/\sqrt{3}}$。
検証:$r_{11}r_{22}r_{33} = \sqrt{2} \cdot \dfrac{\sqrt{6}}{2}
\cdot \dfrac{2}{\sqrt{3}} = \sqrt{\dfrac{2 \cdot 6 \cdot 4}{4 \cdot 3}} = 2$
で、直接計算した $\det A = -2$ の絶対値と一致する。
\end{rei}

\begin{rei}[例題・標準(${}^t\!AA = {}^t\!RR$)]\label{rei:qr-cholesky}
$A = QR$ のとき ${}^t\!AA = {}^t\!RR$ が成り立つことを示し、
例題 \ref{rei:qr-basic3x2} の行列で数値的に確かめよ。\\
\textbf{解.}\ ${}^t\!QQ = E$ を使うだけである:
\[
{}^t\!AA = {}^t\!(QR)(QR) = {}^t\!R\,{}^t\!QQ\,R = {}^t\!RR.
\]
例題 \ref{rei:qr-basic3x2} では
\[
{}^t\!RR = \mat{1 & 0 \\ 2 & 5}\mat{1 & 2 \\ 0 & 5}
= \mat{1 & 2 \\ 2 & 29}, \qquad
{}^t\!AA = \mat{1 & 2 \\ 2 & 29}
\]
で確かに一致する。対称正定値行列 ${}^t\!AA$ が
「三角行列とその転置の積」に書けるというこの形は
\textbf{コレスキー分解}と呼ばれ、数値計算で多用される。
\end{rei}

\begin{rei}[例題・標準(証明:正射影行列は $Q\,{}^t\!Q$)]\label{rei:qr-proj}
列が線形独立な $A = QR$ について、$W = \mathrm{Im}\,A$ への
正射影行列(定理 \ref{thm:lsq-projmat})が
\[
P = A({}^t\!AA)^{-1}{}^t\!A = Q\,{}^t\!Q
\]
と簡単になることを示せ。\\
\textbf{解.}\ 例題 \ref{rei:qr-cholesky} より ${}^t\!AA = {}^t\!RR$ で、
$R$ は正則だから $({}^t\!AA)^{-1} = ({}^t\!RR)^{-1} = R^{-1}({}^t\!R)^{-1}$
(積の逆行列は順序を入れ替えた逆行列の積)。代入して
\[
P = QR \cdot R^{-1}({}^t\!R)^{-1} \cdot {}^t\!R\,{}^t\!Q = Q\,{}^t\!Q.
\]
$Q\,{}^t\!Q = \vect{u}_1{}^t\vect{u}_1 + \cdots + \vect{u}_n{}^t\vect{u}_n$
だから、これは「正規直交基底の各方向への正射影の和」
(定理 \ref{thm:gs-decomp} の公式の行列版)にほかならない。
逆行列の計算が消えてしまうのが QR 分解の御利益である。
\end{rei}

\begin{matome}
\begin{itemize}
\item 列が線形独立な行列は $A = QR$($Q$:列が正規直交系、
$R$:対角成分が正の上三角)と\textbf{ただ一通り}に分解できる
(定理 \ref{thm:qr-main})。存在はグラム・シュミットの
\textbf{計算記録}、一意性は列を左から順に比べれば出る。
\item $r_{jj} = \|\vect{b}_j\|$(直交化の途中で出る「高さ」)、
$r_{ij} = \langle \vect{a}_j, \vect{u}_i \rangle$($i < j$)。
列が従属なら $\vect{b}_j = \vect{0}$ が現れ、この形の分解は存在しない
(反例 \ref{hanrei:qr-dependent})。
\item 最小二乗問題は $R\vect{x} = {}^t\!Q\vect{b}$ の\textbf{後退代入}に
帰着し、数値計算では正規方程式より安定である。
また ${}^t\!AA = {}^t\!RR$、正射影行列は $P = Q\,{}^t\!Q$。
\item 正方行列では $Q$ は直交行列で、
$|\det A| = r_{11}\cdots r_{nn}$(定理 \ref{thm:qr-det}:
「底辺 × 高さ」の体積計算の行列版)。
\item 「対角成分が正」は向きの任意性を固定するための条件で、
外すと一意性が崩れる(反例 \ref{hanrei:qr-sign})。
\end{itemize}
\end{matome}

\section*{章末問題}

\begin{mondai}[基本]\label{q:qr-ex1}
$A = \mat{1 & 1 \\ 1 & -1}$ の QR 分解を求めよ。
$R$ が対角行列になるのはなぜか、一言で説明せよ。
\end{mondai}

\begin{mondai}[基本]\label{q:qr-ex2}
$A = \mat{0 & 2 \\ 3 & 4}$ の QR 分解を求め、
$|\det A| = r_{11}r_{22}$(定理 \ref{thm:qr-det})を確かめよ。
\end{mondai}

\begin{mondai}[基本]\label{q:qr-ex3}
$A = \mat{1 & 1 \\ 0 & 3 \\ 0 & 4}$ の QR 分解を求めよ。
\end{mondai}

\begin{mondai}[基本]\label{q:qr-ex4}
問 \ref{q:qr-ex3} の分解を使い、$\vect{b} = \mat{2 \\ 5 \\ 5}$ に対する
最小二乗解を $R\vect{x} = {}^t\!Q\vect{b}$ で求めよ。
\end{mondai}

\begin{mondai}[標準]\label{q:qr-ex5}
$A = \mat{2 & 1 \\ 1 & 2}$ の QR 分解を求め、
$\det Q = 1$(回転)であることと、
$r_{11}r_{22} = |\det A| = 3$ を確かめよ。
\end{mondai}

\begin{mondai}[標準]\label{q:qr-ex6}
問 \ref{q:lsq} のデータ $(0,0), (1,1), (2,1)$ の当てはめ行列
$A = \mat{1 & 0 \\ 1 & 1 \\ 1 & 2}$, $\vect{b} = \mat{0 \\ 1 \\ 1}$ について
QR 分解を求め、$R\vect{x} = {}^t\!Q\vect{b}$ を解いて、
問 \ref{q:lsq} の答え $a = \dfrac{1}{6}$, $b = \dfrac{1}{2}$ が
再現されることを確かめよ。
\end{mondai}

\begin{mondai}[標準]\label{q:qr-ex7}
問 \ref{q:qr-ex3} の $Q$ の $2$ 本の列に、第 $3$ の列
$\vect{u}_3$ を補って $3$ 次の直交行列を作れ。
\end{mondai}

\begin{mondai}[標準]\label{q:qr-ex8}
問 \ref{q:qr-ex3} の $A$ について、正射影行列 $P = Q\,{}^t\!Q$
(例題 \ref{rei:qr-proj})を計算せよ。さらに
$P\vect{b}$($\vect{b} = {}^t(2, 5, 5)$)が、
問 \ref{q:qr-ex4} の最小二乗解 $\vect{x}^*$ に対する $A\vect{x}^*$ と
一致することを確かめよ。
\end{mondai}

\begin{mondai}[発展]\label{q:qr-ex9}
第 $1$ 列が零でない一般の $2$ 次正方行列
$A = \mat{a & b \\ c & d}$ について、QR 分解の $R$ が
\[
R = \mat{\sqrt{a^2 + c^2} & \dfrac{ab + cd}{\sqrt{a^2 + c^2}} \\[2mm]
0 & \dfrac{|ad - bc|}{\sqrt{a^2 + c^2}}}
\]
となることを示せ($A$ は正則とする)。
例 \ref{rei:qr-2x2} の $A = \mat{3 & 1 \\ 4 & 3}$ で検算せよ。
\end{mondai}

\begin{mondai}[発展]\label{q:qr-ex10}
(QR法の入り口)$A = \mat{2 & 1 \\ 1 & 1}$ を QR 分解し、
積の順序を入れ替えた $A_1 = RQ$ を計算せよ。\\
　(1) $A_1 = {}^t\!QAQ$ を示し、$A_1$ が $A$ と相似
(同じ固有値をもつ)であることを確かめよ
(トレースと行列式で検算せよ)。\\
　(2) $A$ の固有値は $\dfrac{3 \pm \sqrt{5}}{2} \approx 2.618,\ 0.382$
である。$A_1$ の対角成分と見比べ、非対角成分の大きさが
どう変わったかを観察せよ(この操作をくり返すと対角行列に近づいていく
――これが固有値計算の標準アルゴリズム「QR法」である)。
\end{mondai}

\begin{mondai}[証明]\label{q:qr-ex11}
直交行列 $T$ が上三角で、対角成分がすべて正ならば、$T = E$ である
ことを示せ(ヒント:第 $1$ 列から順に。各列はノルム $1$ で、
先に確定した列と直交する)。この事実から定理 \ref{thm:qr-main} の
一意性をもう一度導け($A = QR = Q'R'$ が正方の場合に、
$T = R'R^{-1}$ を考えよ。上三角行列の逆行列と積が上三角であることは
使ってよい)。
\end{mondai}

\begin{mondai}[証明]\label{q:qr-ex12}
(アダマールの不等式)正則な $n$ 次正方行列 $A$ の列を
$\vect{a}_1, \dots, \vect{a}_n$ とすると
\[
|\det A| \le \|\vect{a}_1\| \|\vect{a}_2\| \cdots \|\vect{a}_n\|
\]
が成り立つことを、QR 分解を使って示せ(ヒント:
$\vect{a}_j = r_{1j}\vect{u}_1 + \cdots + r_{jj}\vect{u}_j$ と
ピタゴラスの定理から $r_{jj} \le \|\vect{a}_j\|$)。
等号が成り立つのはどんなときか。幾何的には何を意味するか。
\end{mondai}

\bigskip
{\small
\noindent\textbf{【章末問題の方針】}\quad
まず自力で取り組み、行き詰まったら下の方針を見よ。記述による完全解答は
巻末「問の解答」にある。
\begin{itemize}
\item \textbf{問 \ref{q:qr-ex1}}:$2$ つの列がはじめから直交している。
正規化するだけでよい。
\item \textbf{問 \ref{q:qr-ex2}}:$r_{11} = \|\vect{a}_1\| = 3$。
$\vect{u}_1$ は $\vect{e}_2$ になる。
\item \textbf{問 \ref{q:qr-ex3}}:例題 \ref{rei:qr-basic3x2} と同じ型。
$\vect{a}_1$ はすでに単位ベクトル。
\item \textbf{問 \ref{q:qr-ex4}}:${}^t\!Q\vect{b}$ の成分は
$\langle \vect{b}, \vect{u}_i \rangle$。下の行から解く。
\item \textbf{問 \ref{q:qr-ex5}}:$r_{11} = \sqrt{5}$。分数の根号は
最後に整理する。
\item \textbf{問 \ref{q:qr-ex6}}:$\vect{u}_1 = \dfrac{1}{\sqrt{3}}
{}^t(1,1,1)$、$\vect{b}_2 = {}^t(-1, 0, 1)$。
\item \textbf{問 \ref{q:qr-ex7}}:$\vect{u}_1, \vect{u}_2$ の両方と
直交する単位ベクトルを探す(連立方程式か、成分の観察で)。
\item \textbf{問 \ref{q:qr-ex8}}:$P = \vect{u}_1{}^t\vect{u}_1 +
\vect{u}_2{}^t\vect{u}_2$。$A\vect{x}^*$ は問 \ref{q:qr-ex4} の答えを使う。
\item \textbf{問 \ref{q:qr-ex9}}:定義どおり直交化する。
$\|\vect{b}_2\|$ の計算で $(ab + cd)^2$ を展開し、
$(a^2 + c^2)(b^2 + d^2) - (ab + cd)^2 = (ad - bc)^2$ を使う。
\item \textbf{問 \ref{q:qr-ex10}}:$A = QR$ なら $R = {}^t\!QA$ だから
$RQ = {}^t\!QAQ$。相似なら固有値・トレース・行列式は不変(第III部・第IV部)。
\item \textbf{問 \ref{q:qr-ex11}}:第 $1$ 列は ${}^t(t_{11}, 0, \dots, 0)$ で
ノルム $1$、$t_{11} > 0$ より $t_{11} = 1$。第 $2$ 列は第 $1$ 列と直交…と
順に進む。後半は $Q' \,T = Q$ の形にして、$T$ が直交かつ上三角・対角正に
なることを確かめる。
\item \textbf{問 \ref{q:qr-ex12}}:$\|\vect{a}_j\|^2 = r_{1j}^2 + \cdots
+ r_{jj}^2 \ge r_{jj}^2$。辺々掛けて定理 \ref{thm:qr-det} を使う。
等号は $R$ が対角、すなわち列がもともと直交するとき。
\end{itemize}
}

% =====================================================
\chapter{特異値分解(SVD) ― すべての行列のための分解}\label{chap:svd}
% =====================================================

\section{対角化の限界を超える}

固有値理論には前提があった:行列は\textbf{正方}でなければならず、
美しい結論(直交対角化)は\textbf{対称}行列に限られた。
しかしデータの行列(行 = 標本、列 = 変数)は長方形が当たり前である。
任意の $m \times n$ 行列に通用する分解はないか
――答えが特異値分解である。

\begin{teiri}[特異値分解]
任意の $m \times n$ 行列 $A$(階数 $r$)は、
\[
A = U \Sigma\, {}^t V
\]
と分解できる。ここで $U$ は $m$ 次直交行列、$V$ は $n$ 次直交行列、
$\Sigma$ は対角の位置に
\[
\sigma_1 \ge \sigma_2 \ge \cdots \ge \sigma_r > 0
\]
(\textbf{特異値})を並べ、他を $0$ とした $m \times n$ 行列である。
\end{teiri}

構成は対称行列の理論に帰着する。${}^t\!AA$ は $n$ 次の対称行列で、
しかも ${}^t\vect{x}({}^t\!AA)\vect{x} = \|A\vect{x}\|^2 \ge 0$ だから
固有値はすべて非負である(第V部の言葉で\textbf{半正定値})。そこで
\begin{enumerate}
\item ${}^t\!AA$ を直交対角化し、固有ベクトルを $V$ の列、
固有値の正の平方根 $\sigma_i = \sqrt{\lambda_i}$ を特異値とする。
\item $\vect{u}_i = A\vect{v}_i / \sigma_i$($i \le r$)とすると、
これらは自動的に正規直交系になり、$U$ の列(の最初の $r$ 本)が得られる。
\end{enumerate}

\begin{rei}
$A = \mat{1 & 1 \\ 0 & 0}$。
${}^t\!AA = \mat{1 & 1 \\ 1 & 1}$ の固有値は $2, 0$、固有ベクトルは
$\vect{v}_1 = \dfrac{1}{\sqrt{2}}\mat{1 \\ 1}$,
$\vect{v}_2 = \dfrac{1}{\sqrt{2}}\mat{1 \\ -1}$。
特異値は $\sigma_1 = \sqrt{2}$(階数 $1$ なので一つだけ)。
$\vect{u}_1 = A\vect{v}_1/\sigma_1 = \mat{1 \\ 0}$、これに直交する
$\vect{u}_2 = \mat{0 \\ 1}$ を補えば
\[
A = \mat{1 & 0 \\ 0 & 1}\mat{\sqrt{2} & 0 \\ 0 & 0}
\cdot \frac{1}{\sqrt{2}}\mat{1 & 1 \\ 1 & -1}
\]
($最後の行列は {}^tV$)。展開して元に戻ることを確かめてほしい。
\end{rei}

\section{幾何的な意味 ― すべての線形写像は「回して、伸ばして、回す」}

$A = U\Sigma{}^tV$ を変換として読むと:
${}^tV$ は直交行列(回転・鏡映)、$\Sigma$ は座標軸ごとの伸縮、
$U$ も直交行列。すなわち

\begin{itembox}[l]{SVD の幾何}
\textbf{どんな線形写像も「回転 → 軸に沿った伸縮 → 回転」の
三段階に分解できる}。線形写像は単位球面を必ず楕円体にうつし、
その楕円体の主軸の長さが特異値 $\sigma_1, \dots, \sigma_r$ である。
\end{itembox}

固有値が「向きを変えない方向の伸縮率」だったのに対し、
特異値は「入力の直交軸($\vect{v}_i$)と出力の直交軸($\vect{u}_i$)を
別々に選んでよいときの伸縮率」である。入力と出力の空間が違ってもよく、
正方でなくてもよいのはこのためである。
最大特異値 $\sigma_1$ は「$A$ がベクトルを最大何倍に引き伸ばせるか」を表し、
行列の「大きさ」の標準的な尺度になる。

なお、$A$ が対称で固有値が非負なら、SVD は直交対角化(第V部)と一致する。
SVD は\textbf{スペクトル定理の、全行列への拡張}なのである。

\begin{mondai}\label{q:svd}
$A = \mat{0 & 2 \\ 1 & 0}$ の特異値分解を求めよ。
\end{mondai}

% =====================================================
\chapter{低ランク近似と擬似逆行列}\label{chap:pinv}
% =====================================================

\section{低ランク近似 ― データ圧縮の数学}

SVD は和の形にも書ける(スペクトル分解の類似):
\[
A = \sigma_1 \vect{u}_1{}^t\vect{v}_1 + \sigma_2 \vect{u}_2{}^t\vect{v}_2
+ \cdots + \sigma_r \vect{u}_r{}^t\vect{v}_r.
\]
各項 $\vect{u}_i{}^t\vect{v}_i$ は階数 $1$ の行列であり、
$A$ は「階数 $1$ の部品 $r$ 個の、重み $\sigma_i$ 付きの重ね合わせ」である。
特異値は大きい順に並んでいるから、\textbf{上位 $k$ 項で打ち切れば、
$A$ のよい近似が得られる}はずである。実際、次が成り立つ。

\begin{teiri}[エッカート・ヤングの定理]
上位 $k$ 項の和 $A_k = \sigma_1\vect{u}_1{}^t\vect{v}_1 + \cdots +
\sigma_k\vect{u}_k{}^t\vect{v}_k$ は、
\textbf{階数 $k$ 以下のすべての行列の中で $A$ に最も近い}。
\end{teiri}

\begin{itembox}[l]{深掘り:画像圧縮・主成分分析・推薦システム}
$1000 \times 1000$ 画素のモノクロ画像は $100$ 万個の数からなる行列だが、
自然画像の特異値は急速に減衰するため、上位 $k = 50$ 項だけ、すなわち
$\sigma_i, \vect{u}_i, \vect{v}_i$ の約 $50 \times 2001 \approx 10$ 万個の数
(元の $1$ 割)で、見た目にほぼ劣化のない近似が得られる。
これが SVD による画像圧縮の原理である。

統計学の\textbf{主成分分析(PCA)}も同じ数学である。データ行列に SVD を施すと、
$\vect{v}_1$ は「データが最も大きくばらつく方向」(第 $1$ 主成分)になる
(分散共分散行列――第V部の深掘りに登場した対称行列――の
固有ベクトルと一致する)。高次元データを少数の主成分で要約する操作は、
エッカート・ヤングの意味での最良低ランク近似にほかならない。

さらに、推薦システム(ユーザー × 商品の巨大な評価行列から
未知の評価を予測する問題)も、「評価行列は本質的に低ランクである
(好みは少数の潜在因子で決まる)」という仮説のもとで
低ランク近似を行う技術であり、Netflix の推薦コンテストで
一躍有名になった。\textbf{特異値の減衰は、データに潜む
「本質的な次元の低さ」の表現}なのである。
\end{itembox}

\section{擬似逆行列 ― すべての方程式に「公式の答え」を}

最後に、本書第I部からの長い物語に決着をつける。
$A\vect{x} = \vect{b}$ の解は「一意・無数・なし」の三択だった。
擬似逆行列は、\textbf{三つの場合すべてに、ただ一つの標準的な答え}を与える。

\begin{teigi}[ムーア・ペンローズの擬似逆行列]
$A = U\Sigma{}^tV$ に対し、$\Sigma$ の零でない対角成分を逆数にして
転置した行列を $\Sigma^{+}$ とし、
\[
A^{+} = V\,\Sigma^{+}\,{}^t U
\]
を $A$ の\textbf{擬似逆行列}という($A$ が正則なら $A^{+} = A^{-1}$ に一致する)。
\end{teigi}

\begin{teiri}
$\vect{x}^{*} = A^{+}\vect{b}$ は、
\begin{enumerate}
\item 解が一意に存在するなら、その解そのもの。
\item 解が存在しないなら、残差 $\|\vect{b} - A\vect{x}\|$ を最小にする
\textbf{最小二乗解}。
\item 解(または最小二乗解)が無数にあるなら、その中で
ノルム $\|\vect{x}\|$ が\textbf{最小}のもの。
\end{enumerate}
\end{teiri}

「解けないなら最も近くへ、答えが多すぎるなら最も慎ましいものを」
――第I部で三つに割れた世界が、一つの式 $\vect{x}^* = A^{+}\vect{b}$ に
統一された。実際の数値計算ライブラリで連立方程式や回帰を解くとき、
内部で動いているのはまさにこの SVD 経由の擬似逆行列である。

\begin{mondai}\label{q:pinv}
$A = \mat{1 & 1 \\ 0 & 0}$(本章の例で SVD 済み)について、\\
　(1) $A^{+}$ を求めよ。\\
　(2) 解の存在しない方程式 $A\vect{x} = \mat{1 \\ 1}$ に対し
$\vect{x}^{*} = A^{+}\mat{1 \\ 1}$ を計算し、これが
最小二乗解の中でノルム最小のものになっていることを
(最小二乗解の全体を求めて)確かめよ。
\end{mondai}
